\subsection{拓扑向量空间中凸集的分离}

定义 1. --- 实拓扑向量空间 E 的两个非空集 A, B 称为被一个闭超平面 H 分离，如果 A 含于 H 所确定的一个闭半空间中，而 B 含于另一个闭半空间中。

定义 2. --- 实拓扑向量空间的两个非空集 A, B 称为被闭超平面 H 严格分离，如果 A 含于 H 所确定的一个开半空间中，而 B 含于另一个开半空间中。

命题 1. --- 设 A 是一个非空开凸集，B 是实拓扑向量空间 E 中的一个非空凸集；若 A 不与 B 相交，则存在一个把 A 与 B 分离的闭超平面。

因为集 \(\mathbf{C} = \mathbf{A} - \mathbf{B}\) 是开的、凸的 (II, 第 9 页, prop. 7) 且非空，又有 \(0 \notin \mathbf{C}\) 。由 II, 第 36 页的定理 1，存在一个连续线性型 \(f \neq 0\) （\(\mathbf{E}\) 上的），使 \(f(z) > 0\) （在 \(\mathbf{C}\) 中） 。于是，对所有 \(x \in \mathbf{A}\) 及 \(y \in \mathbf{B}\) ，有 \(f(x) > f(y)\) 。记 \(\alpha = \inf_{x \in \mathbf{A}} f(x)\) ；\(\alpha\) 是有限的，且 \(f(x) \geqslant \alpha\) （对所有 \(x \in \mathbf{A}\) ），\(f(y) \leqslant \alpha\) （对所有 \(y \in \mathbf{B}\) ） ；闭超平面 \(\mathbf{H}\) （方程为 \(f(z) = \alpha\) ）把 \(\mathbf{A}\) 与 \(\mathbf{B}\) 分离。

注记. --- 1) 超平面 H 不与 A 相交 (II, 第 15 页, prop. 1；若 A 和 B 是两个互不相交的非空开凸集，则存在一个把 A 与 B 严格分离的闭超平面。

2) 然而，当 B 不是开集时，未必存在把 A 与 B 严格分离的闭超平面，即使 E 是有限维的，甚至即使 \(\overline{A}\) 不与 \(\overline{B}\) 相交 (II, 第 78 页, exerc. 12)。

定义 3. --- 对向量空间 E 的一个子集 A，超平面 H 称为 A 的支撑超平面，如果 H 至少包含 A 的一点，且 A 的所有点都位于 H 的同一侧。

设 f 是 E 上一个不恒为零的线性型；说方程为 \(f(x) = \alpha\) 的超平面是 A 的一个支撑超平面，意味着 \(\alpha\) 是集 \(f(A) \subset R\) 的最小元或最大元。换言之，存在 A 的一个平行于方程 \(f(x) = 0\) 的超平面的支撑超平面，必须且只需集 \(f(A)\) 的两个界之一有限且属于 \(f(A)\) 。

命题 2. --- 设 A 是拓扑向量空间 E 的一个非空紧子集。对 E 中每个闭超平面 H，存在 A 的一个平行于 H 的支撑超平面。

因为，若 \(f(x) = \gamma\) 是 H 的一个方程，其中 f 是 E 中的一个连续线性型，则 f 在 A 上的限制连续，因而有界且在 A 上达到其界 (GT, IV, § 6.1, th. 1)。

这表明存在一个或两个平行于 H 的 A 的支撑超平面；第一种情形只能在 A 完全含于一个平行于 H 的超平面时出现。

命题 3. --- 在拓扑向量空间 E 中，设 A 是一个内部非空的闭凸集。则 \(\mathbf{A}\) 的每个支撑超平面都是闭的，且 \(\mathbf{A}\) 的每个边界点都属于 \(\mathbf{A}\) 的至少一个支撑超平面。

A 的每个支撑超平面都是闭的，因为 A 的所有点都位于超平面的同一侧 (II, 第 15 页, prop. 17)。又若 \(x_0\) 是 A 的一个边界点，则 \(x_0\) 不属于非空开凸集 \(\mathring{\mathrm{A}}\) ；由 II, 第 36 页的 th. 1，存在一个包含 \(x_0\) 且不与 \(\mathring{\mathrm{A}}\) 相交的超平面 H。由于 A 是 \(\mathring{\mathrm{A}}\) 的闭包 (II, 第 14 页, prop. 16 的 cor. 1)，由 II, 第 15 页的 prop. 17 推出 H 是 A 的一个支撑超平面。

